In this work, we introduced a new method for learning SDEs from population-level snapshot data. Our approach is based on matching state--time distributions using a least squares scheme in a distributional space. Our proposed method handles real-life challenges such as unknown and state-dependent volatility, missing dimensions, and diagnostics for performance. Overall, our experiments indicate that our proposed framework outperforms existing methods in a wide range of applications. While many scientific applications (including those above) feature useful models, it can nonetheless be a limitation that our method requires specification of an appropriate model family. A poorly chosen model family can lead to bad performance. An overly wide family can face identifiability challenges. In fact, even the complete time series of marginal distributions need not always uniquely determine the drift and volatility functions of the SDE; see \cref{app:identifiability} for further discussion. Finally, our method is not simulation-free, so it requires compute that scales with sample size and state dimension; development of a simulation-free method is an interesting future direction.

\textbf{Ethics statement.} 
This work does not involve human subjects or personally identifiable information, and the datasets we use (PBMC, Repressilator simulations, Lotka–Volterra, Gulf of Mexico HYCOM reanalysis) are all either publicly available or simulated. We believe the main applications of our method are in scientific forecasting tasks, such as gene regulation, ecological modeling, and environmental monitoring, which we view as positive in their social impact. Nonetheless, we note that any forecasting tool could in principle be misapplied in sensitive contexts (e.g., healthcare or policy). We adhere to the ICLR Code of Ethics in the conduct and reporting of this work.  

\textbf{Reproducibility statement.} We have taken several steps to ensure reproducibility. 
All theoretical results are stated in the main text (\cref{sec:method}) with complete proofs in \cref{app:proof-main-prop}. 
Experimental setups, including SDE parameterizations, training details, and evaluation metrics, are described in \cref{sec:experiments} as well as in \cref{app:further-exp-details}, where we have one appendix for each experiment. Code to reproduce all experiments is provided at \url{https://anonymous.4open.science/r/snapMMD-DD84/}. 
